\begin{quote}
%This paper provides causal evidence on the effects of beat policing on public security and police effectiveness. Using rich administrative and survey data from Chile, we study \emph{Plan Cuadrante}, a reform that shifted policing toward a beat-policing model by subdividing urban municipalities into smaller geographic units---i.e., quadrants---and permanently assigning each officer to one quadrant rather than rotating officers across the municipality. Exploiting the staggered implementation of the program across municipalities, we find large improvements across multiple outcomes. Among surveyed households, twelve-month victimization rates declined by 10 percentage points (36\%), and administrative records show a 10\% decline in crimes reported to the police. The reform also improved security perceptions. The share of households perceiving increased criminal activity fell by 15 percentage points (36\%), and, consistent with this pattern, household investment in private security measures declined by 7.7 percentage points (37\%). In some municipalities, implementation of \emph{Plan Cuadrante} was accompanied by significant increases in police resources. However, our mechanism analyses indicate that our findings are not driven by resource increases alone. The beat-policing components of the reform---i.e., permanent officer assignment and greater local knowledge---appear to be key to reducing crime and improving security perceptions.
This paper provides causal evidence on the effects of beat policing on public security and police effectiveness. Using rich administrative and survey data from Chile, we study \emph{Plan Cuadrante}, a reform that shifted policing toward a beat-policing model by subdividing urban municipalities into smaller geographic units---i.e., quadrants---and permanently assigning officers to these quadrants rather than rotating officers across the municipality. Exploiting the staggered implementation of the program across municipalities, we find large improvements in both public safety and security perceptions. Among surveyed households, twelve-month victimization rates declined by 10 percentage points (36\%), and administrative records show a 10\% decline in crimes reported to the police. Arrests increased by 14\% even as crime declined. The reform also improved security perceptions. The share of households perceiving increased criminal activity fell by 15 percentage points (36\%), and, consistent with this pattern, household investment in private security measures declined by 7.7 percentage points (37\%). In some municipalities, implementation of \emph{Plan Cuadrante} was accompanied by significant increases in police resources. However, our mechanism analyses indicate that our findings are not driven by resource increases alone, as they remain sizeable in municipalities with resource increases similar to those in control municipalities. The beat-policing components of the reform---i.e., permanent officer assignment and greater local knowledge---appear to be key to reducing crime and improving security perceptions.
\\[12pt]
\noindent\textbf{Keywords:} police strategy, beat policing, territorial deployment.\\
%\textbf{JEL codes:} 

\end{quote}